\documentclass[10pt,a4paper]{article}
\usepackage[T1]{fontenc}
\usepackage[utf8]{inputenc}
\usepackage{lmodern}
\usepackage[english]{babel}
\usepackage[margin=1.8cm]{geometry}
\usepackage{graphicx,booktabs,tabularx,enumitem,microtype,xcolor}
\usepackage[hidelinks]{hyperref}
\hypersetup{pdftitle={Tandem MCP Architecture 4.1.0},pdfauthor={Alessandro Nocentini (ALENOC)}}
\setlist{nosep,leftmargin=1.2em}
\setlength{\parindent}{0pt}
\setlength{\parskip}{4pt}
\newcommand{\code}[1]{\texttt{\detokenize{#1}}}

\title{\vspace{-1.5em}Tandem MCP Architecture\\\large Version 4.1.0: cloud orchestrator, local executor}
\author{Alessandro Nocentini (ALENOC)\\\small\url{https://github.com/ALENOC/tandem-mcp}}
\date{November 2026}

\begin{document}
\maketitle
\vspace{-1em}

\section{Overview}
Tandem MCP is a Model Context Protocol server that lets the model in your chat hand part of a coding
task to a model running on your own machine. The cloud orchestrator (Claude Code, Codex, Antigravity,
Gemini CLI or another MCP client) keeps design, contracts and review. The local model (LM Studio, Ollama,
llama.cpp, vLLM or any OpenAI-compatible server) reads large files and writes code inside the limits of
the contract, and Tandem checks every result before it reaches the orchestrator.

\begin{center}
\includegraphics[width=0.62\textwidth]{diagram-en.pdf}
\end{center}

\section{Core modules}
\begin{tabularx}{\textwidth}{@{}lX@{}}
\toprule
Module & Responsibility\\
\midrule
\code{server.py} & 20 MCP tools (4 in the autonomous profile), strict contract validation, receipts, secret redaction\\
\code{engine.py} & Server discovery, capability probe, role routing (digest and delegate), retry with backoff\\
\code{digest.py} & Chunked map-reduce reading, lexical quote checks, ungrounded numbers, schema merge\\
\code{agent.py} & Prefix-stable executor loop, lazy write compaction, multi-edit \code{replace_in_file}, journal, batches, rollback\\
\code{sandbox.py} & Glob matching, command parser, outbound symlink refusal, SHA-256 stamp, Bubblewrap jail\\
\code{ui_test.py} & \code{tandem:ui_test} scenarios from the contract, run in headless Chromium\\
\code{web_check.py} & Built-in runtime, animation and 40 FPS frame-rate gate for web deliverables\\
\code{router.py} & \code{auto_route}: deterministic tools, direct answer or local delegation\\
\code{host_runner.py} & Host-managed wait for a job, with no cloud polling\\
\code{store.py} & Persistence under \code{TANDEM_HOME}, cross-process locks, startup reconciliation\\
\bottomrule
\end{tabularx}

\section{Delegation and the autonomous profile}
\begin{itemize}
\item \textbf{Contract:} \code{delegate} needs a goal, \code{write_globs}, \code{verify_commands} that fail
  until the work is done, and atomic constraints, one requirement each. In the autonomous profile the
  server splits compound constraints instead of refusing the contract.
\item \textbf{Baseline and protection:} checks run before any edit; test files are read-only and restored if touched.
\item \textbf{Local loop and repair:} the conversation only grows at its end so the model server reuses its
  cached prompt; replies allow up to 12,000 tokens (\code{TANDEM_REPLY_TOKENS}). On \code{finish}, Tandem checks
  syntax, runs shell checks in the sandbox (sending slow passing checks over \code{TANDEM_SLOW_CHECK_S} back once),
  runs \code{tandem:ui_test} and the 40 FPS \code{web_check} in headless Chromium, and verifies cited lines.
  Trimmed outputs keep failing lines; three identical failures in a row stop early as \code{stuck}.
\item \textbf{Autonomous profile} (\code{TANDEM_PROFILE=autonomous}): four tools. \code{delegate} blocks while
  the loop runs locally, with no cloud call, bounded by \code{TANDEM_JOB_SECONDS} (3600),
  \code{TANDEM_MAX_REPAIRS} (20) and \code{TANDEM_MAX_STEPS} (160). A progress notification every 30 seconds
  keeps the client from aborting the call. The orchestrator gets one \code{ACCEPTED} line, or only the
  checks that still fail.
\end{itemize}

\section{Security model and rollback}
\begin{itemize}
\item \textbf{Sandbox profiles} (\code{TANDEM_SANDBOX}): \code{secure} (fail-closed) by default on Linux,
  \code{compatible} on macOS and Windows, \code{trusted} disables Bubblewrap.
\item \textbf{Bubblewrap jail:} system folders and the Python runtime are mounted read-only and the
  workspace is the only writable folder; host \code{$HOME}, \code{/tmp} and \code{/run} are empty; PID and IPC
  namespaces are separate, and so is the network unless the task allows it.
\item \textbf{Command parser:} every part of a pipeline must match \code{allow_run} and not the deny list;
  command substitution and heredocs are refused.
\item \textbf{Backup and rollback:} original bytes are journaled per job (8 MiB per file, 256 MiB per job);
  an incomplete rollback sets the state \code{partial_rollback}.
\item \textbf{Secret redaction:} tool outputs mask private keys, tokens and secret environment assignments.
\end{itemize}

\section{Measured savings}
With Opus 5.5 as the orchestrator and the autonomous profile (4.1.0 series, four medium and large tasks, four
counterbalanced pairs each), Tandem cut total cloud spend by 57\% and cloud cost per successful result by 61\%
(\$0.217 vs \$0.556), with 67\% fewer output tokens (excluding two final-run infrastructure incidents: 10/15
successes each, -61\% total and per success, -72\% output tokens). Median task time dropped from 25 min to
19 min (40--45\% faster per local call), and the 40 FPS gate lifted 3D portal frame rates from 4--20 FPS to
38--60 FPS. Small edits under 30 lines remain direct (+10\% in 4.0.0). Local model: Qwen3.8 27B IQ3\_S on an
AMD Radeon RX 6800 series GPU (16 GB). Details: \code{docs/BENCHMARKS.md} and
\url{https://alenoc.github.io/tandem-mcp/}.

\end{document}
